% Options for packages loaded elsewhere
\PassOptionsToPackage{unicode}{hyperref}
\PassOptionsToPackage{hyphens}{url}
\PassOptionsToPackage{table}{xcolor}
%
\documentclass[
  11pt,
]{report}
\usepackage{amsmath,amssymb}
\usepackage{iftex}
\ifPDFTeX
  \usepackage[T1]{fontenc}
  \usepackage[utf8]{inputenc}
  \usepackage{textcomp} % provide euro and other symbols
\else % if luatex or xetex
  % pdfLaTeX-only repository policy: the original unicode-math/fontspec
  % branch is unreachable under pdfLaTeX and is replaced with a guard.
  \PackageError{cpython-suite}{Compile with pdfLaTeX only}{}
\fi
\usepackage{lmodern}
\ifPDFTeX\else
  % xetex/luatex font selection
\fi
% Use upquote if available, for straight quotes in verbatim environments
\IfFileExists{upquote.sty}{\usepackage{upquote}}{}
\IfFileExists{microtype.sty}{% use microtype if available
  \usepackage[]{microtype}
  \UseMicrotypeSet[protrusion]{basicmath} % disable protrusion for tt fonts
}{}
\makeatletter
\@ifundefined{KOMAClassName}{% if non-KOMA class
  \IfFileExists{parskip.sty}{%
    \usepackage{parskip}
  }{% else
    \setlength{\parindent}{0pt}
    \setlength{\parskip}{6pt plus 2pt minus 1pt}}
}{% if KOMA class
  \KOMAoptions{parskip=half}}
\makeatother
\usepackage{xcolor}
\usepackage{listings}
\newcommand{\passthrough}[1]{#1}
\lstset{defaultdialect=[5.3]Lua}
\lstset{defaultdialect=[x86masm]Assembler}
\usepackage{longtable,booktabs,array}
\usepackage{calc} % for calculating minipage widths
% Correct order of tables after \paragraph or \subparagraph
\usepackage{etoolbox}
\makeatletter
\patchcmd\longtable{\par}{\if@noskipsec\mbox{}\fi\par}{}{}
\makeatother
% Allow footnotes in longtable head/foot
\IfFileExists{footnotehyper.sty}{\usepackage{footnotehyper}}{\usepackage{footnote}}
\makesavenoteenv{longtable}
\setlength{\emergencystretch}{3em} % prevent overfull lines
\providecommand{\tightlist}{%
  \setlength{\itemsep}{0pt}\setlength{\parskip}{0pt}}
\setcounter{secnumdepth}{5}
% CPython Engineering Suite shared presentation layer -- pdfLaTeX compatible
\usepackage[T1]{fontenc}
\usepackage{lmodern}
\usepackage{microtype}
\usepackage{geometry}
\geometry{letterpaper,margin=0.91in,headheight=16pt,headsep=19pt}
\usepackage{xcolor} % [table] passed at top of file (option-clash fix)
\definecolor{CPNavy}{HTML}{13283B}
\definecolor{CPBlue}{HTML}{176BA0}
\definecolor{CPGold}{HTML}{D8A73C}
\definecolor{CPMuted}{HTML}{566979}
\definecolor{CPBg}{HTML}{EFF5F8}
\usepackage{titlesec}
\titleformat{\chapter}[display]{\normalfont\sffamily\bfseries\color{CPNavy}}{\large\color{CPBlue}CHAPTER \thechapter}{9pt}{\Huge}[\vspace{6pt}\color{CPGold}\titlerule]
\titleformat{\section}{\Large\bfseries\sffamily\color{CPBlue}}{\thesection}{0.7em}{}
\titleformat{\subsection}{\large\bfseries\sffamily\color{CPNavy}}{\thesubsection}{0.7em}{}
\usepackage{fancyhdr}
\pagestyle{fancy}
\fancyhf{}
\fancyhead[L]{\footnotesize\color{CPMuted}\sffamily CPYTHON ENGINEERING SUITE}
\fancyhead[R]{\footnotesize\color{CPMuted}\sffamily 3.16.0a1+dev SNAPSHOT}
\fancyfoot[C]{\footnotesize\color{CPMuted}\thepage}
\renewcommand{\headrulewidth}{0.3pt}
\AtBeginDocument{\urlstyle{tt}}
\lstset{basicstyle=\ttfamily\small,columns=fullflexible,keepspaces=true,breaklines=true,breakatwhitespace=false,showstringspaces=false,frame=single,framesep=5pt,rulecolor=\color{CPBg},backgroundcolor=\color{CPBg},xleftmargin=3pt,xrightmargin=3pt}
\usepackage{tcolorbox}
\tcbuselibrary{breakable,skins}
\newtcolorbox{sourcebox}[1][]{enhanced,breakable,sharp corners,boxrule=0.4pt,colback=CPBg,colframe=CPBlue,borderline west={3pt}{0pt}{CPBlue},left=10pt,right=8pt,top=6pt,bottom=6pt,#1}
\newtcolorbox{cautionbox}[1][]{enhanced,breakable,sharp corners,boxrule=0.4pt,colback=CPBg,colframe=CPGold,borderline west={3pt}{0pt}{CPGold},left=10pt,right=8pt,top=6pt,bottom=6pt,#1}
\usepackage{tikz}
\usetikzlibrary{arrows.meta,positioning,shapes.geometric,calc}
\usepackage{enumitem}
\setlist{itemsep=2pt,topsep=4pt}
\usepackage{booktabs}
\usepackage{longtable}
\AtBeginEnvironment{longtable}{\small}
\usepackage{array}
\setlength{\parskip}{5pt}
\setlength{\parindent}{0pt}
\setcounter{tocdepth}{2}
\usepackage{hyperref}
\hypersetup{colorlinks=true,linkcolor=CPBlue,urlcolor=CPBlue,citecolor=CPBlue,pdfauthor={CPython Engineering Documentation Suite}}
\makeatletter
\renewcommand{\maketitle}{%
  \begin{titlepage}\pagecolor{CPNavy}\color{white}\sffamily
  \vspace*{0.8in}
  {\color{CPGold}\rule{2.1in}{4pt}}\\[16pt]
  {\Large\bfseries CPYTHON ENGINEERING DOCUMENTATION SUITE}\\[24pt]
  {\Huge\bfseries \@title\par}\vspace{19pt}
  \vfill
  {\normalsize Source baseline: CPython 3.16.0 alpha 1+dev repository snapshot}\par
  {\normalsize Prepared 16 October 2026}\par
  {\small\color{white!70} Source-traced engineering commentary; not an official Python specification.}\vspace{0.65in}
  \end{titlepage}\nopagecolor\color{black}
  }
\makeatother
\ifLuaTeX
  \usepackage{selnolig}  % disable illegal ligatures
\fi
\usepackage{bookmark}
\IfFileExists{xurl.sty}{\usepackage{xurl}}{} % add URL line breaks if available
\urlstyle{same}
\hypersetup{
  pdftitle={CPython Interpreter Architecture},
  hidelinks,
  pdfcreator={LaTeX via pandoc}}

\title{CPython Interpreter Architecture}
\usepackage{etoolbox}
\makeatletter
\providecommand{\subtitle}[1]{% add subtitle to \maketitle
  \apptocmd{\@title}{\par {\large #1 \par}}{}{}
}
\makeatother
\subtitle{Volume 03 \textbar{} Runtime, evaluation, frames,
specialization, JIT and lifecycle}
\author{}
\date{}

\begin{document}
\maketitle

{
\setcounter{tocdepth}{1}
\tableofcontents
}
\chapter{Architecture baseline and interpretation
limits}\label{architecture-baseline-and-interpretation-limits}

\begin{sourcebox}
\textbf{Evidence classification.} \textbf{SOURCE} = statements supported by material present in the uploaded snapshot; \textbf{IMPLEMENTATION} = CPython-specific internals, not guaranteed Python semantics; \textbf{ENGINEERING} = original recommendations, examples, checks, and exercises. The source baseline is \textbf{3.16.0 alpha 1+dev}; missing repository file bodies are noted as evidence limitations.
\end{sourcebox}

This volume presents the internals described in
\path{InternalDocs/}, which explicitly characterizes
itself as documentation \textbf{for CPython maintainers}. These
mechanisms are not guaranteed by the Python language specification.
\path{InternalDocs/interpreter.md} identifies
\path{Python/ceval.c} as the bytecode interpreter's
entry point and says the instruction implementations are generated from
\path{Python/bytecodes.c} by a case generator. Those
paths occur in the repository tree but their complete bodies were not
included in the supplied text export; this guide therefore treats the
internal documentation as the direct supporting evidence.

\textbf{Baseline:} \passthrough{\lstinline!README.rst!} identifies
3.16.0 alpha 1+dev. An interpreter implementation document written for
that baseline can describe the design, but cannot establish a stable
ABI, a permanent opcode set, or final-release behavior.

\chapter{Repository anatomy and dependency
direction}\label{repository-anatomy-and-dependency-direction}

\path{InternalDocs/structure.md} describes a useful
four-way mapping: user docs, public-library implementation, native
implementation and tests. In the core runtime,
\passthrough{\lstinline!Include/!} provides public headers,
\path{Include/internal/} internal declarations,
\passthrough{\lstinline!Objects/!} object implementations,
\passthrough{\lstinline!Python/!} execution and runtime machinery,
\passthrough{\lstinline!Parser/!} syntax processing,
\passthrough{\lstinline!Modules/!} extension and built-in modules,
\passthrough{\lstinline!Lib/!} Python implementations and tests, and
\passthrough{\lstinline!Tools/!} generators and developer tooling.

\begin{longtable}[]{@{}
  >{\raggedright\arraybackslash}p{(\columnwidth - 4\tabcolsep) * \real{0.3333}}
  >{\raggedright\arraybackslash}p{(\columnwidth - 4\tabcolsep) * \real{0.3333}}
  >{\raggedright\arraybackslash}p{(\columnwidth - 4\tabcolsep) * \real{0.3333}}@{}}
\toprule\noalign{}
\begin{minipage}[b]{\linewidth}\raggedright
Layer
\end{minipage} & \begin{minipage}[b]{\linewidth}\raggedright
Representative source references
\end{minipage} & \begin{minipage}[b]{\linewidth}\raggedright
Nature of contract
\end{minipage} \\
\midrule\noalign{}
\endhead
\bottomrule\noalign{}
\endlastfoot
Language & \path{Doc/reference/executionmodel.rst} &
Python program behavior \\
Runtime public C surface & \path{Include/Python.h},
\passthrough{\lstinline!Doc/c-api/!} & Versioned C API and ABI
policies \\
Object implementation & \passthrough{\lstinline!Objects/!},
\path{Include/cpython/} & CPython-specific
implementation \\
Interpreter execution &
\path{InternalDocs/interpreter.md},
\path{Python/ceval.h} & Internal, not stable \\
Compiler & \path{InternalDocs/compiler.md},
\path{Python/assemble.c} & Internal representation \\
Diagnostics and tests & \passthrough{\lstinline!Lib/test/!},
\passthrough{\lstinline!Tools/gdb/!},
\path{InternalDocs/} & Verification and
maintainability \\
\end{longtable}

\textbf{Interpretation rule:} a file existing under
\path{Include/internal/} or containing
\passthrough{\lstinline!\_Py!} does not mean third-party extensions may
depend on it. Consult \path{Doc/c-api/stable.rst} for
the API stability hierarchy.

\chapter{Interpreter lifecycle and state
ownership}\label{interpreter-lifecycle-and-state-ownership}

The runtime creates interpreter state and per-thread execution state
before executing Python code. Frames provide the dynamic context of a
running code object. Runtime initialization and finalization are
different from module initialization: a host embedding Python must
manage initialization, configuration, thread state and shutdown as one
lifecycle. For extension code, multi-phase module initialization and
per-module state are preferred whenever isolation across interpreters is
required.

\textbf{Evidence:} \path{Doc/c-api/init.rst},
\path{Doc/c-api/interp-lifecycle.rst},
\path{Doc/c-api/subinterpreters.rst},
\path{InternalDocs/frames.md}, and
\path{Include/cpython/initconfig.h}.

\textbf{ENGINEERING state model:}

\begin{enumerate}
\def\labelenumi{\arabic{enumi}.}
\tightlist
\item
  Establish runtime configuration (filesystem paths, environment and
  security policy).
\item
  Initialize the runtime and relevant interpreter(s).
\item
  Load modules with explicit state ownership.
\item
  Execute code while preserving exception, stack and thread invariants.
\item
  Release module resources, tear down interpreters and finalize safely.
\end{enumerate}

Do not infer a single global interpreter, one thread, or shared module
state in designs intended for subinterpreters or free-threaded builds.

\chapter{Code objects, frames and
execution}\label{code-objects-frames-and-execution}

\section{Static code information versus dynamic execution
state}\label{static-code-information-versus-dynamic-execution-state}

\path{InternalDocs/code_objects.md} describes code
objects as products of compilation that store bytecode and static
metadata. \path{InternalDocs/frames.md} distinguishes
frame representations and their lifetime. An instruction execution uses
an active frame and its evaluation stack; an object representing a
user-visible Python frame is not necessarily allocated for every
internal activation in every implementation scenario.

\textbf{Pipeline model:}

\begin{lstlisting}
Python source -> tokenizer -> PEG parser -> AST
             -> instruction sequence -> optimized CFG
             -> assembled code object
             -> frame activation -> eval/dispatch -> Python-visible result
\end{lstlisting}

The compiled code object typically contains instructions, constants,
names and information used for exception handling. Bytecode is an
\textbf{implementation format}, not a standardized serialization
protocol. Do not persist \passthrough{\lstinline!.pyc!} files or opcode
sequences as a stable cross-version contract.

\section{Dispatch, inline caches and
specialization}\label{dispatch-inline-caches-and-specialization}

\path{InternalDocs/interpreter.md} explains the
adaptive interpreter and instruction specialization. Dispatching
bytecode executes operations against the current stack and frame. Inline
caches store feedback needed for specialized execution paths, and
deoptimization must preserve Python semantics. Debug output from
\passthrough{\lstinline!dis!} is useful, but bytecode details can evolve
even between successive CPython builds.

\begin{lstlisting}[language=Python]
import dis

def score(record):
    return record["points"] * 2

for _ in range(1000):
    score({"points": 7})

dis.dis(score)
\end{lstlisting}

\textbf{ENGINEERING:} Treat disassembly as observational evidence about
a pinned interpreter build; never assume a particular instruction
sequence is permanent or derive source-level correctness from the
specialization path chosen.

\chapter{Stack references, exception propagation and
suspension}\label{stack-references-exception-propagation-and-suspension}

\section{\texorpdfstring{\texttt{\_PyStackRef}
ownership}{\_PyStackRef ownership}}\label{pystackref-ownership}

\path{InternalDocs/stackrefs.md} describes tagged
pointer-sized stack references carrying metadata such as ownership
information and tagged small integers. This internal representation is a
performance and lifetime-management tool: it is not equivalent to a
public \passthrough{\lstinline!PyObject*!} and is not an extension ABI.
Operations for duplication, borrowing, closing and clearing have
distinct ownership implications. Stack reference handling must preserve
lifetime even during exceptional control flow.

\section{Exception tables and
unwinding}\label{exception-tables-and-unwinding}

\path{InternalDocs/exception_handling.md} describes
the bytecode-level machinery that routes exceptions to handlers and
manages unwinding. User code observes
\passthrough{\lstinline!try!}/\passthrough{\lstinline!except!}/\passthrough{\lstinline!finally!}
semantics; internal handlers may be encoded through exception metadata
rather than older opcode shapes. Generator and coroutine suspension
preserve execution state until resumed or finalized; see
\path{InternalDocs/generators.md} and
\path{InternalDocs/frames.md}.

\textbf{Correctness probe:} run the same function through success,
exception, and cancellation paths. If an optimization changes a stack
reference's closing behavior, examine every exit path, not only the
successful return.

\chapter{Adaptive tier and experimental
JIT}\label{adaptive-tier-and-experimental-jit}

\path{InternalDocs/jit.md} describes the adaptive
interpreter and an optional trace-based JIT, historically called tiers 1
and 2. The trace recorder can translate sequences into micro-operations
and optimize across instruction boundaries. The document describes
hotness detection involving instructions such as
\path{JUMP_BACKWARD} and
\passthrough{\lstinline!RESUME!} and an initialization route to tracing.
This description is specific to a JIT-capable CPython build; it does
\textbf{not} imply all Python installations use a JIT.

\begin{longtable}[]{@{}
  >{\raggedright\arraybackslash}p{(\columnwidth - 4\tabcolsep) * \real{0.3333}}
  >{\raggedright\arraybackslash}p{(\columnwidth - 4\tabcolsep) * \real{0.3333}}
  >{\raggedright\arraybackslash}p{(\columnwidth - 4\tabcolsep) * \real{0.3333}}@{}}
\toprule\noalign{}
\begin{minipage}[b]{\linewidth}\raggedright
Component
\end{minipage} & \begin{minipage}[b]{\linewidth}\raggedright
Interpreter role
\end{minipage} & \begin{minipage}[b]{\linewidth}\raggedright
Diagnostic question
\end{minipage} \\
\midrule\noalign{}
\endhead
\bottomrule\noalign{}
\endlastfoot
Code object & Stores compiled instructions and metadata & Which source
block produced this code? \\
Frame & Holds dynamic execution state & Which values and exception state
are live? \\
Adaptive interpreter & Executes/specializes instructions & Is
specialization preserving the same result? \\
Trace recorder & Captures hot execution regions on supported builds &
What guard or exit ends the trace? \\
JIT execution & Applies optimization across sequences & Does
deoptimization return to correct state? \\
\end{longtable}

\textbf{Implementation caution:} A JIT is an optional and evolving
implementation technique. For performance experiments, record executable
version, compilation settings, process environment, warm-up policy,
workload, CPU, memory and confidence intervals. A timing difference does
not prove a code-quality improvement without profiling and variance
analysis.

\chapter{Observability and operational
diagnostics}\label{observability-and-operational-diagnostics}

CPython provides multiple visibility layers: tracebacks,
\passthrough{\lstinline!inspect!}, \passthrough{\lstinline!dis!},
\path{sys.monitoring},
\passthrough{\lstinline!tracemalloc!}, CPU profiling and GDB helpers.
Each has a different measurement impact.
\path{InternalDocs/profiling_binary_format.md}
discusses internal profiler format details; treat such artifacts as
version-pinned unless the documentation specifically promises stability.

\begin{lstlisting}[language=Python]
import sys
import platform

print(sys.version)
print(sys.implementation.name)
print(platform.platform())
print(sys.getswitchinterval())
\end{lstlisting}

For incident investigations, capture reproducible inputs and interpreter
build metadata first. Distinguish high-level symptoms (exception,
deadlock, unexpectedly high allocation rate) from low-level hypotheses
(specialization regression, interpreter thread state, allocator
contention).

\chapter{Architectural risks and review
controls}\label{architectural-risks-and-review-controls}

\begin{itemize}
\tightlist
\item
  \textbf{State leakage:} global/static extension state may violate
  subinterpreter isolation.
\item
  \textbf{Ownership ambiguity:} borrowed, owned and stack references
  cannot be interchanged casually.
\item
  \textbf{Runtime assumptions:} a free-threaded build changes
  concurrency behavior, not the programmer's obligation to synchronize
  shared mutation.
\item
  \textbf{Unstable internals:} using private
  \passthrough{\lstinline!\_Py!} APIs creates source and binary
  compatibility risk.
\item
  \textbf{Measurement distortion:} tracing and profiling alter costs;
  compare instrumented and uninstrumented baselines.
\item
  \textbf{Import-time side effects:} imports occur in
  runtime/interpreter/module contexts and can accidentally bind system
  resources.
\end{itemize}

\chapter{Hands-on interpreter architecture
laboratories}\label{hands-on-interpreter-architecture-laboratories}

\section{Lab 1: Trace a function from source to
frame}\label{lab-1-trace-a-function-from-source-to-frame}

Write \passthrough{\lstinline!def total(values): return sum(values)!}.
Use \passthrough{\lstinline!compile!}, inspect
\passthrough{\lstinline!\_\_code\_\_!} and
\passthrough{\lstinline!dis.dis!}, and record constants, names, local
variable metadata and disassembly. Cross-check behavior with
\path{InternalDocs/code_objects.md} and
\path{InternalDocs/interpreter.md}.

\section{Lab 2: Inspect exception
recovery}\label{lab-2-inspect-exception-recovery}

Create nested \passthrough{\lstinline!try!} /
\passthrough{\lstinline!finally!} code that intentionally raises.
Compare the result with the displayed code-object exception table and
the Python-visible traceback. Explain which obligations are semantic
guarantees and which table contents are implementation-specific.

\section{Lab 3: Detect optimization-sensitive measurement
errors}\label{lab-3-detect-optimization-sensitive-measurement-errors}

Benchmark a warmed and unwarmed function repeatedly; run with clearly
recorded interpreter flags. Report distribution of measurements, not the
single fastest trial, and discuss possible interpreter
specialization/JIT confounders.

\section{Lab 4: Module state isolation
review}\label{lab-4-module-state-isolation-review}

Choose an extension-module design with static global counters. Redesign
it around per-module state. Provide a threat or correctness narrative
for two interpreter instances and concurrent execution.

\textbf{Deliverables:} short architecture diagram, source-path evidence
log, executable probes, and a two-column summary of language guarantees
versus CPython implementation observations.

\chapter{Source concordance}\label{source-concordance}

\textbf{Source bodies inspected:}
\path{InternalDocs/README.md},
\passthrough{\lstinline!structure.md!},
\path{interpreter.md},
\path{code_objects.md},
\passthrough{\lstinline!frames.md!},
\passthrough{\lstinline!stackrefs.md!},
\path{exception_handling.md},
\path{generators.md},
\passthrough{\lstinline!jit.md!},
\path{profiling_binary_format.md};
\path{Doc/c-api/init.rst} and
\path{Doc/c-api/stable.rst}.

\textbf{Referenced but not fully supplied as bodies:} key
generated/interpreter sources such as
\path{Python/ceval.c},
\path{Python/bytecodes.c}, and
\path{Python/optimizer.c}, referenced by the internal
explanations. Detailed per-function claims about those absent source
bodies have intentionally been avoided.

\end{document}
